\documentclass[11pt,a4paper]{article}
\usepackage[margin=23mm]{geometry}
\usepackage{amsmath,amssymb,array}
\input{glyphtounicode}
\pdfgentounicode=1
\usepackage[hidelinks]{hyperref}
\hypersetup{hypertexnames=false}
\setlength{\parindent}{0pt}
\setlength{\parskip}{5pt}
\newcommand{\pos}[1]{\left[#1\right]_+}
\newcommand{\Reflect}{\mathcal R}
\begin{document}
\begin{center}
{\Large P-H: a conditional recovery-work point node}\\[4pt]
15 October 2026 --- proposed exact law, analytical qualification only
\end{center}

\textbf{Decision status.} P-H's causal slow-entry/recovery direction is
user-approved at CP-D1. The exact reduction below is proposed at CP-D2,
not yet approved or calibrated. It substitutes the compulsory point-service
component of the unchanged H1/T1 predecessor. No simulator, traffic rollout,
gradient computation, fitting or optimisation was performed.

\section*{1. Shared inputs and the single-component delta}

Retain the predecessor's uniform footprint $[0,\ell]$, preparation length
$L_p\geq\ell$, terminal length $b>0$, $L=L_p+b$, no subsequent external
arrivals, initially empty terminal, continuing symmetry $a_0=b_0$, total
mass $M=\ell(e_0+a_0+b_0)$, and one advice $v\in[v_-,v_+]$.
Use the same positive wave speed $w$, jam density $k_j$, and triangular family
\begin{equation}
 c(s)=\frac{wk_js}{w+s},\quad k_c(s)=\frac{wk_j}{w+s},\quad
 r=e_0+a_0,\quad v_+r<R\leq c(v_-),\quad
 v_+b_0\leq R_B\leq c(v_-).
 \tag{1}
\end{equation}
$R$ is the fixed ordinary shared service ceiling. It is not a prescribed
future receiving history. H1's paired early transfer remains
\begin{equation}
 \dot e=-\frac{v}{h}\min\left\{e,\frac{k_c(v)-r+e}{\gamma}\right\},
 \qquad a=r-e,\qquad b_0\text{ constant},\quad h>0,\ \gamma\geq1.
 \tag{2}
\end{equation}
During $t_f^p=(L_p-\ell)/v<t<t_r^p=L_p/v$, terminal entrance rates
are $(\dot F_E,\dot F_A,\dot F_B)=v(e(t),a(t),b_0)$, and zero otherwise.
Define histories, extended by zero to negative times, by
\begin{equation}
 P_j(t)=F_j(t-b/v),\quad N_j(t)=F_j(t)-Y_j(t),\quad
 q_j(t)=P_j(t)-Y_j(t),\qquad j=E,A,B.
 \tag{3}
\end{equation}
$N$ is physical stock; $q$ is eligible unserved mass, not spatial queue
length. All geometry/parameters are shared across advice values and observers.

\textbf{Delta.} Remove the predecessor's instantaneous extra E work factor
$\alpha$ and the identity $Y_A=P_A$. Retain ordinary one-for-one mass
admission, strict A priority and independent B service. Replace the node
budget by the causal recovery history below. H1's early work factor
$\gamma$, physical opportunity scale $h$, cutoff and control are unchanged.
There is no new lateral opportunity in the terminal or spatial recovery zone.

\newpage
\section*{2. An isolated restriction with a finite recovery duration}

\textbf{Source basis versus deduction.} The point case in
\href{https://repository.tudelft.nl/file/File_c424c1b3-7374-436b-a6a7-abdf69063aa6}{Leclercq et al. (2016), section 2, eq. (1), printed p4}
uses bounded acceleration and backward waves to calculate cumulative
passing capacity. The recovery envelope below is our isolated-case
deduction, capped at the retained ceiling $R$. The many-entry convolution
on the next page is a proposed closure, not that source's native law.

Let $m_{\rm ref}=1$ vehicle-equivalent fix physical mass normalization.
It is not a packet size or an integer admission. Define the speed at which
an isolated restriction becomes neutral to $R$:
\begin{equation}
 u_R=\frac{wR}{wk_j-R},\qquad c(u_R)=R,\qquad u_R\leq v_-.
 \tag{4}
\end{equation}
Propose a representative insertion speed $0\leq u_*<u_R$ and positive
recovery acceleration $a_{\rm r}>0$. These are effective shared parameters;
they are not an Eulerian exit-speed field inferred from (3). An isolated
prototype at age $\xi$ has
\begin{equation}
 x(\xi)=u_*\xi+\tfrac12a_{\rm r}\xi^2,\quad
 \theta=\xi+x(\xi)/w,\quad
 U(\theta)=\sqrt{(w+u_*)^2+2wa_{\rm r}\theta}-w.
 \tag{5}
\end{equation}
$\theta$ is age as seen by the boundary after a backward wave, whereas
$\xi$ is prototype travel age. Equivalently, differentiating the isolated
cumulative expression $wk_j(\theta-\xi(\theta))$ gives $c(U(\theta))$.
Only until the restriction reaches $R$ is that expression used.

The resulting loss profile is
\begin{align}
 \Theta&=\frac{u_R-u_*}{a_{\rm r}}
       +\frac{u_R^2-u_*^2}{2a_{\rm r}w},\nonumber\\
 D(\theta)&=
 \begin{cases} R-c(U(\theta)),&0\leq\theta\leq\Theta,\\
                0,&\theta>\Theta.
 \end{cases}
 \tag{6}
\end{align}
$D$ is nonnegative, decreasing, bounded by $D_0=R-c(u_*)$, continuous
at $\Theta$ and Lipschitz on $[0,\infty)$; $\|D'\|_1=D_0$.
In particular $u_*=0$ gives $D_0=R$, with no singularity or positive-speed
epsilon. Its integrated missing work is exactly
\begin{equation}
 W=\int_0^\Theta D(\theta)d\theta
   =\frac{k_c(u_R)(u_R-u_*)^2}{2a_{\rm r}}<\infty.
 \tag{7}
\end{equation}
For example, use $d\theta/du=(w+u)/(a_{\rm r}w)$ and
$[R-c(u)]d\theta=k_c(u_R)(u_R-u)du/a_{\rm r}$ to verify (7).
Because $R,w,k_j,u_*,a_{\rm r}$ are fixed and $v\geq u_R$, this profile
is fixed across advices. Advice still changes (2), arrivals and admissions;
no additional direct advice-to-insertion-speed relation is asserted.

\newpage
\section*{3. Proposed continuous history and conservative priority service}

\textbf{Constitutive assumption.} Completed continuous E mass leaves
mass-weighted copies of (6). Their loss demands add; the finite budget
saturates at zero if those demands exhaust it:
\begin{equation}
 p(t)=\frac{1}{m_{\rm ref}}\int_0^tD(t-s)\,dY_E(s),\qquad
 \mu(t)=\pos{R-p(t)},\qquad p(0)=0.
 \tag{8}
\end{equation}
Only $Y_E$ in the last $\Theta$ seconds is needed. It is predicted history,
not initial hidden information, a list of individual cars, or stored
longitudinal trajectories. Linear overlap of loss demands is an explicit
reduction of physical interactions, not an exact wave/void calculation.
Nonnegative residual service in (8) is part of this proposed service law,
not an unapproved numerical clipping or recovery smoothing.

For $j=E,A$ define eligible demand almost everywhere by
\begin{equation}
 d_j(t)=\begin{cases}c(v),&q_j(t)>0,\\
                          \dot P_j(t),&q_j(t)=0.
          \end{cases}\qquad
 g_A=\min\{d_A,\mu\},\quad
 g_E=\min\{d_E,\pos{\mu-g_A}\},\quad \dot Y_j=g_j.
 \tag{9}
\end{equation}
$Y_B=P_B$ on the retained independent B domain. At an empty queue,
available arrivals restrict departure; a nonempty queue can send at $c(v)$.
Since $0\leq\mu\leq R\leq c(v)$, a nonempty combined queue receives all
available shared service. A has priority but can be delayed by (8).
Exactly $g_jdt$ mass leaves, never a work-weighted amount. There is no
forward escape for E or second compulsory merge of early-transferred mass.

\begin{center}
\begin{tabular}{p{0.38\linewidth}p{0.5\linewidth}}
Quantity & Units and role\\ \hline
$Y,F,P,N,q,m_{\rm ref},W$ & Vehicle-equivalent mass (history, stock or integrated work).\\
$R,c,D,p,\mu,g,d$ & Vehicle-equivalents per second (service or demand).\\
$u_*,u_R,v,w$ & Length per second.\\
$a_{\rm r},\Theta$ & Length per second squared; seconds, respectively.\\
$p$ & Service-loss demand, not physical traffic stock.\\
\end{tabular}
\end{center}

For finite completed E mass $M_E$,
\begin{equation}
 \int_0^\infty p(t)dt=\frac{M_E}{m_{\rm ref}}W.
 \tag{10}
\end{equation}
This is Tonelli's identity for the nonnegative finite-support kernel,
conditional at first on the eventual admitted mass. Even before clearance,
the left side is bounded above by $MW/m_{\rm ref}$ because $Y_E\leq P_E\leq M$.
An E burst may have completed while $p$ remains positive. Its influence
expires by $\Theta$ after its last admission; its vehicle-time does not
continue in the local objective merely because $p$ persists.

\newpage
\section*{4. Existence and uniqueness for this particular node}

The queue-indicator form (9) need not be used to infer uniqueness.
For a continuous candidate cumulative service $K(t)=\int_0^t\mu(s)ds$,
define one-sided reflection, with histories starting at zero, by
\begin{equation}
 \Reflect(X)(t)=X(t)-\min\left\{0,\inf_{0\leq s\leq t}X(s)\right\}.
 \tag{11}
\end{equation}
For fixed $K$, the work-conserving strict-priority solution is
\begin{align}
 q_A&=\Reflect(P_A-K),\qquad Q=\Reflect(P_A+P_E-K),\nonumber\\
 q_E&=Q-q_A,\qquad
 Y_A=P_A-q_A,\qquad Y_E=P_E+q_A-Q.
 \tag{12}
\end{align}
Nondecreasing $P_E$ gives $Q\geq q_A$. Checking the cases $q_A>0$,
$q_A=0<Q$ and $Q=0$ recovers (9); each departure is nonnegative,
nondecreasing in cumulative mass and bounded by eligible arrivals.
The total rate is at most $\dot K\leq R$, hence also respects each
lane's sending ceiling. All histories are absolutely continuous.

Reflection is 2-Lipschitz in the uniform norm. Holding arrivals fixed,
\begin{equation}
 \|\Delta Y_E\|_\infty\leq4\|\Delta K\|_\infty.
 \tag{13}
\end{equation}
Integration by parts, using $Y_E(0)=0$ and the absolutely continuous $D$
on nonnegative ages, gives
\begin{equation}
 p(t)=\frac{D_0Y_E(t)+\int_0^tD'(t-s)Y_E(s)ds}{m_{\rm ref}},\qquad
 \|\Delta p\|_\infty\leq\frac{8D_0}{m_{\rm ref}}\|\Delta K\|_\infty.
 \tag{14}
\end{equation}
Thus the causal map
\begin{equation}
 K\ \longmapsto\ \int_0^t\pos{R-p[Y_E[K]](s)}ds
 \tag{15}
\end{equation}
is a contraction on any interval of length
$\eta<m_{\rm ref}/(8D_0)$ with fixed identical past histories. The closed
space consists of cumulative $K$ starting at zero with slopes in $[0,R]$.
The map preserves it. Successive intervals therefore give a unique global
causal solution; the same estimates and Gronwall give continuous dependence
on arrivals and the kernel on finite time intervals. If $D_0=0$, $K=Rt$
directly. The interval length is a proof device, not a simulation time step.

No future insertion schedule or receiving trace is supplied. Continuous
admissions generate (8), which determines subsequent admissions uniquely.
This is a proof for (6)--(9), not an appeal to a general invariant-node
theorem for an arbitrary history-dependent junction.

\newpage
\section*{5. Physical storage and finite full-cohort clearance}

Retain the sufficient all-action domain
\begin{equation}
 M< H=k_jb,\qquad 0\leq\dot F_j\leq c(v),\qquad F_j\leq M.
 \tag{16}
\end{equation}
The accepted empty-link vacancy is $k_jwt-F_j(t)$ until $b/w$, then
$Y_j(t-b/w)+H-F_j(t)$. Before $b/w$, it is positive because
$c(v)<k_jw$; afterwards it is at least $H-M>0$. This proof needs only
nonnegative departures, not immediate A service. H1 therefore keeps its
unrestricted entrance/handoff on (16) even when either terminal queue
stops temporarily. Physical stock $F_j-Y_j$ is conserved, bounded and
never replaced by the service-memory variable. The domain protects H1;
it does not demonstrate general spillback into the preparation window.

All traffic is eligible by $t_*=L/v$. B is then clear. Suppose combined
eligible A/E mass $Q(t_*)>0$ has not yet cleared by a time $T>t_*$.
With no subsequent arrivals, $Q$ is nonincreasing and remains positive
throughout that interval. Work conservation and (8) give
\begin{align}
 \dot Y_A+\dot Y_E&=\mu,\qquad R\leq\mu+p,\nonumber\\
 R(T-t_*)&\leq\int_{t_*}^T\mu(t)dt+\int_{t_*}^Tp(t)dt\nonumber\\
 &\leq Q(t_*)+\frac{\overline M_E}{m_{\rm ref}}W,
 \qquad \overline M_E=P_E(t_*).
 \tag{17}
\end{align}
The last inequality uses only $Y_E\leq\overline M_E$, not prior clearance.
Letting $T$ grow contradicts (17), proving finite clearance. In particular
\begin{equation}
 T_{\rm clear}\leq \frac{L}{v}
   +\frac{Q(L/v)+(\overline M_E/m_{\rm ref})W}{R}.
 \tag{18}
\end{equation}
If $Q(t_*)=0$, the cohort is already clear. A common, conservative horizon
for all admitted advices and same-$M$ observer comparisons is
\begin{equation}
 \overline T=\frac{L}{v_-}+\frac{M+(M/m_{\rm ref})W}{R}.
 \tag{19}
\end{equation}
All physical and eligible queues are empty after actual clearance. Any
remaining service influence is not counted as an unfinished vehicle.

\textbf{Meaning.} No positive service floor is inserted: $p\geq R$ can
give $\mu=0$. Finite total admitted mass creates a finite integrated burden,
so temporary closure cannot persist forever. Equations (17)--(19) replace
the predecessor's inapplicable positive instantaneous service bound.

\newpage
\section*{6. MP6: a nonempty conditional family with A delay}

This symbolic existence construction establishes possibility, not a
calibrated regime, numerical scenario, or preferred advice. It can require
slow recovery and a long terminal. Fix one advice $v=v_+$ satisfying (1),
let $e_0,a_0=b_0>0$, and choose $\gamma>1$ so
\begin{equation}
 \delta=k_c(v)-r>0,\qquad (\gamma-1)e_0>\delta.
 \tag{20}
\end{equation}
The receiving branch of (2) is strictly active initially. Let $L_p=\ell$,
so the front enters the terminal at time zero. Measured from the first
eligible arrival, at $b/v$, the arrival rates are
\begin{equation}
 \lambda_A(t)=v(r-e(t))\nearrow\lambda=vr<R,\qquad
 \lambda_E(t)=ve(t)>0,\quad \lambda_A+\lambda_E=\lambda
 \quad (0<t<\ell/v).
 \tag{21}
\end{equation}
The formal extension of (2) to all positive times has $e\downarrow0$,
finite positive $\int_0^\infty e(t)dt$, and that integral scales with $h$.
Choose $h$ so the potential E arrivals in this extension exceed
$m_{\rm ref}(R-\lambda)/D_0$.

First use $a_{\rm r}=0$ only as an \emph{auxiliary finite-time limit}:
$D(t)=D_0>0$ and $\mu=\pos{R-D_0Y_E/m_{\rm ref}}$. It is not an
admitted recovery law or a clearance example. While $\mu>\lambda$, all
arrivals leave and $Y_E$ accumulates the potential E arrivals. The chosen
integral forces $\mu$ to reach $\lambda$ at a finite time when
$\lambda_E>0$. A positive E departure rate makes $\mu<\lambda$ just
afterwards. From then on $\mu$ cannot increase in this auxiliary limit.
Choose a time $t_1$ with $\mu(t_1)<\lambda$. Since $\lambda_A\to\lambda$,
eventually A demand exceeds $\mu(t_1)$ by a positive amount. If A has not
already queued, its eligible stock must then start increasing. Thus
$q_A(t_2)>0$ at some finite $t_2$.

As $a_{\rm r}\downarrow0$, the positive-recovery kernels (6) converge
uniformly to $D_0$ on every fixed finite interval. The finite-interval
continuous dependence from (12)--(15) preserves the strict inequality
$q_A(t_2)>0$ for sufficiently small \emph{positive} $a_{\rm r}$. Choose
$\ell/v>t_2$; the finite footprint then reproduces that prefix exactly.
Finally choose finite $b>M/k_j$. This changes the common eligibility delay,
not the translated node argument. Initial symmetry and B's bounds can be
retained, and (16)--(19) apply to the admitted positive-recovery law.

Hence the \emph{same finite cohort} can have active early receiving
limitation, positive continuing-lane waiting caused by forced-entry history,
protected H1 storage and finite clearance. This establishes no empirical
acceleration range, natural queue size, action split or interior optimum.
The zero-recovery auxiliary limit is used only to construct a finite prefix;
it is explicitly excluded from the admitted model parameter set.

\newpage
\section*{7. Objective, information and sensitivity consequences}

From the predecessor's uniform footprint and unchanged paired early
conservation, local complete cohort time is
\begin{align}
 J(v)&=\int_0^{\overline T}\left[M-Y_E(t)-Y_A(t)-Y_B(t)\right]dt\nonumber\\
 &=\frac{M(L-\ell/2)}{v}
   +\int_0^{\overline T}(q_E(t)+q_A(t))dt.
 \tag{22}
\end{align}
B's queue is zero on (1). This includes all pre-line waiting and the full
clearance tail. The service record in (8) is not another vehicle queue.
Post-line vehicle-time remains outside $J$ even if its influence remains.
The fixed bound (19) removes artificial horizon truncation; it does not
hold physical arrival or queue events fixed as advice changes.

Under the retained symmetry, aggregate initial density and ending density
recover $a_0=b_0=(\rho_0-e_0)/2$. Shared parameters and empty initial memory
then fix (2)--(3); uniqueness above fixes (8)--(9). Therefore A+E and G
have exactly the same objective function for every admitted advice.
At points of differentiability they also have the same derivative.
This is model-relative sufficiency, not a minimum-information theorem for
SUMO or a discrete-order necessity result. Unknown initial service memory
would reopen G2 and is not introduced here.

The causal sensitivity path includes
\begin{equation}
 v\ \longrightarrow\ \hbox{early transfer and arrival clocks}
 \ \longrightarrow\ Y_E\ \longrightarrow\ p
 \ \longrightarrow\ (Y_A,Y_E)\ \longrightarrow\ J.
 \tag{23}
\end{equation}
The fixed prototype profile has no direct $v$ dependence in this member.
Admission times, priorities, active constraints and the recovery influence
of those admissions do depend on advice. A derivative that freezes them,
omits A waiting, or stops before clearance is not the complete derivative.
Uniqueness and Lipschitz dependence alone do not prove differentiability
at every switching event or AD/finite-difference correctness.

\textbf{What is not inherited.} The predecessor's positive action certificate
at large $\alpha$, its common-fastest parameter range, and its derivative
formulas apply to its unchanged instantaneous-work member only. The old
aggregate-history loading construction can still initialise the same H1
family with a newly known empty node, but a new control-information advantage
requires its own objective comparison. No such advantage, positive interior
optimum or useful gradient has been proved or computed here.

\newpage
\section*{8. Limiting checks, physical limits and the gate}

\textbf{No E.} If $e_0=0$, then $Y_E=p=0$, $\mu=R$ and (1) admits A
without waiting. The remaining complete travel time decreases with advice.

\textbf{Negligible recovery burden.} As $u_*\uparrow u_R$ or
$a_{\rm r}\to\infty$, $W\to0$. The limiting node has ordinary shared
service $R$. It does not recover the old extra-work member with $\alpha>1$.

\textbf{Stationary connection to a simpler rival.} If E departures have
been constant at rate $g_E$ over a full recovery window and the node is
saturated with $p<R$, then
\begin{equation}
 p=\frac{W}{m_{\rm ref}}g_E,\qquad
 g_A+\left(1+\frac{W}{m_{\rm ref}}\right)g_E=R.
 \tag{24}
\end{equation}
This is our stationary consequence, not the source many-event capacity
formula. It explains why instantaneous work can be a serious average-flow
rival. Finite-burst timing, rather than novelty of a work factor, must
justify the retained history.

\textbf{A delay is permitted, not inevitable.} For constant A arrivals
$\lambda_A<R$, empty initial memory, decreasing $D$ and strict A priority,
the barrier $p\leq R-\lambda_A$ is preserved: when service reaches
$\lambda_A$, E receives zero service and only recovery can change $p$,
so it cannot rise through that barrier. A stays unqueued in this fluid
case. Increasing A demand in (21) avoids that special protection. The model
does not predict that every slow insertion produces a backward shockwave.

\textbf{Main limits.} A representative fixed insertion speed is a closure
parameter, not a traffic-state-derived field. Additive loss demands and
budget saturation approximate overlapping physical restrictions. No source
guarantees their finite-burst fidelity. Full downstream trajectories,
wave/void interactions, stochastic vehicle heterogeneity and post-line
vehicle-time are not reconstructed by (8). The conservative storage domain
deliberately protects H1 from terminal feedback. These limits must appear
in any paper claim, parameter qualification and SUMO comparison.

\textbf{Gate.} Confirming CP-D2 would admit this exact law only as a
conditional elementary primary model. Independent development evidence,
renewed action/sensitivity analysis and the accepted SUMO S-P then S-L
qualification remain. Strong aggregate-history inference, A+E, scalar
search, analytical/finite-difference gradients and the phenomenological
capacity-drop node are retained. No empirical law, new observation, numerical
parameter set, executable evaluator or next roadmap stage is approved here.

\textbf{Reference record.} The accompanying
\href{../../extraction/research_readiness/vsl_ph_source_and_adaptation_audit.md}{source audit}
identifies all four papers rechecked in this pass and their access limits,
including the 2011 conference version and Jin's continuous-LTM and
capacity-drop interfaces. All new proofs above are for the stated proposed
member. The accepted H1/T1 source/PDF are preserved unchanged.
\end{document}
